% !TEX root = ../../main.tex
% F5 - Abstract (two parts). Paragraphs 1-3 translate the Vietnamese abstract (Part I, figures from report-data-guide.md);
% paragraph 4 condenses the abstract of the research report (Part II).
\chapter*{Abstract}
\addcontentsline{toc}{chapter}{Abstract}

When a person must be found in the data of a surveillance camera system, operators usually review the video by hand, which takes a long time and easily misses things as the number of cameras and the recording time grow. This project builds a web application that helps search for people in camera data, letting the user describe the person in one of three forms: a reference image, an English description, or a set of predefined appearance attributes. The system only proposes and ranks results; the decision whether a result is the person being sought remains with the user.

The application has two flows. The indexing flow runs in a background worker: video from an RTSP stream or an uploaded file is sampled, people are detected with YOLO11n and tracked with ByteTrack to group the frames of one person into an appearance, and the representative frame is encoded into a vector with the RaSa model. The data is written to three stores: PostgreSQL for the business data, Milvus for the vectors and MinIO for the frames. The search flow runs in a Flask application server: the query is encoded into the same vector space and matched only within the cameras the user is allowed to see. A React web application serves three roles, the Administrator, the Operator and the Viewer, with functions for managing cameras and AI models, searching, viewing results in their original frame and saving results into case files.

The application was deployed on a CPU-only laptop and tested with the seven cameras of the WILDTRACK dataset, of which 1,488 tracks from the first 120 seconds of each video were indexed without errors. It passes 740 unit tests together with integration and end-to-end tests, and the demonstration script for all three roles ran automatically twice in a row with 15/15 steps passed. Search by image performs well, while search by text and by attributes by the vectors alone is nearly ineffective, mainly because the model's accuracy lies in its cross-modal re-ranker; adding that re-ranker over the first 128 candidates raises Recall@8 of text search from 0 to 0.423 at about 17 seconds per query on the CPU. The analysis is about seven times slower than real time, so the system processes the cameras in turn.

To address the weakness of text search, the second part of the project studies a semantic reasoning layer on top of CLIP retrieval. Each person crop is described offline by a vision-language model (Qwen2.5-VL-3B) into eighteen human-readable attributes with per-region image-quality statistics; at query time, CLIP returns a candidate pool and a multi-agent funnel driven by a small language model (Gemini Flash-Lite) re-ranks it from the attribute evidence, over a coarse and a fine round. On 100 identities of a 496-identity subset of Market-1501, the four-stage funnel raises R@10 from 33.0\% for CLIP to 56.0\%, rescuing 28 of the 67 identities CLIP missed while losing 5 of the 33 it had found. Together, the two parts show that a complete multimodal person search system can be built from existing models on commodity hardware, and that reasoning over readable attribute evidence is a promising way to improve its weakest form of search.
